% ===== CV TEMPLATE: RF / Microwave / Industrial NPI =====
% Emphasis: RF/microwave engineering depth, system industrialisation, hardware
%           development lifecycle. Research credentials kept but secondary.
% Claude Code instruction: copy to cv/cv.tex. For RF-focused roles, emphasise
% the RF/microwave bullet points. For NPI/industrialisation roles, expand the
% production engineering section. Do not alter facts or dates.
\documentclass[a4paper,11pt]{article}
\usepackage[utf8]{inputenc}
\usepackage[T1]{fontenc}
\usepackage[margin=1in]{geometry}
\usepackage{setspace}
\usepackage{hyperref}
\usepackage{enumitem}
\setstretch{1.0}
\pagestyle{empty}

\begin{document}

%====================
% Header
%====================
\begin{center}
    {\LARGE \textbf{Bo Yuan}} \\
    \vspace{0.15cm}
    +46 769\,612\,787 \enspace $\cdot$ \enspace
    \href{mailto:yuanbo330lab@gmail.com}{yuanbo330lab@gmail.com} \enspace $\cdot$ \enspace
    \href{mailto:boyua@kth.se}{boyua@kth.se} \\
    Stockholm / Uppsala, Sweden
\end{center}

\vspace{0.3cm}

%====================
% Objective
%====================
\section*{Objective}

Engineering professional with a strong foundation in RF and microwave systems,
complex electromechanical hardware development, and full-lifecycle system
industrialisation. Five years as CTO leading the design and productisation of a
high-power microwave plasma system from first principles to small-series production.
Currently completing an MSc in Electromagnetics and Plasma Engineering at KTH.
Seeking technically demanding roles in RF/microwave engineering or hardware
industrialisation in the Stockholm / Uppsala area.

%====================
% Education
%====================
\section*{Education}

\noindent\textbf{KTH Royal Institute of Technology}, Stockholm, Sweden
\hfill 09/2024 -- 07/2026 (expected) \\
\textit{M.Sc.\ in Electromagnetics, Fusion and Space Engineering --- Plasma Track} \\
Relevant coursework: Electromagnetic Waves, Plasma Physics, Microwave Modelling,
Plasma--Surface Interaction, Vacuum Technology, Fusion Systems

\medskip
\noindent\textbf{Xi'an University of Technology}, China \hfill 09/2013 -- 07/2018 \\
\textit{B.Eng.\ in Electrical Engineering and Automation}

%====================
% Technical Skills
%====================
\section*{Technical Skills}

\begin{itemize}[noitemsep, topsep=2pt, leftmargin=*]
    \item \textbf{RF \& Microwave Engineering:}
          High-power microwave systems (2.45~GHz, 10~kW magnetron-based),
          waveguide structures, impedance matching, stub tuners, resonant cavity
          design (TM$_{012}$ mode), mode conversion, vacuum-compatible RF components,
          sealing window design, microwave sliding shorts

    \item \textbf{Simulation Tools:}
          ANSYS HFSS ($\sim$2\,yr), COMSOL Multiphysics ($\sim$3\,yr),
          Keysight ADS ($\sim$1\,yr)

    \item \textbf{System \& Industrialisation:}
          System architecture, modular design, design for manufacturing (DFM/DFA),
          NPI, production line setup, manufacturing routing, assembly sequence design,
          incoming inspection, module-level testing, SOP and work instruction authoring

    \item \textbf{Electronics \& Control:}
          PCB design and layout (Altium Designer, $\sim$10\,yr, multi-layer mixed-signal),
          PCBA design coordination and EMS supplier interface,
          Siemens S7-1200 PLC integration, control architecture definition

    \item \textbf{Mechanical \& Vacuum:}
          SolidWorks ($\sim$10\,yr), high-vacuum system design,
          vacuum-tight sealing, leak detection, CNC fixture design

    \item \textbf{Software \& Data:}
          Python (engineering scripts, data analysis), C++ (basic), Odoo ERP
\end{itemize}

%====================
% Professional Experience
%====================
\section*{Professional Experience}

\noindent\textbf{Sichuan 330 Semiconductor Ltd.}, China
\hfill 07/2019 -- 06/2024 \\
\textit{Chief Technology Officer / Co-Founder}

\noindent Led the full lifecycle development and industrialisation of a high-power
MPCVD (Microwave Plasma Chemical Vapor Deposition) system --- 2.45~GHz, 10~kW
magnetron-based, high-vacuum --- from first-principles design through prototyping,
productisation, and small-series production, delivering $\sim$20 complete systems.
Company subsequently acquired.

\medskip
\noindent\textit{RF and Microwave Engineering:}
\begin{itemize}[noitemsep, topsep=2pt, leftmargin=*]
    \item Designed a TM$_{012}$-mode coaxial-tuned resonant reactor from first
          principles using HFSS and COMSOL; optimised EM field distribution for
          uniform plasma generation; enlarged effective deposition area to 75\,mm
    \item Replaced lossy TE--TEM mode converters with low-loss stub-tuned coaxial
          transitions, reducing heat generation and improving power transfer efficiency
    \item Designed vacuum-compatible waveguide sealing windows for reliable
          high-power microwave transmission under vacuum
    \item Designed adjustable microwave sliding shorts with choke structures for
          impedance tuning with reduced RF leakage
\end{itemize}

\medskip
\noindent\textit{System Architecture and Industrialisation:}
\begin{itemize}[noitemsep, topsep=2pt, leftmargin=*]
    \item Defined overall system architecture and modular design strategy;
          decoupled subsystems to enable parallel manufacturing and independent testing
    \item Drove DFM decisions: consolidated multi-part assemblies into CNC-machined
          monolithic components to reduce tolerance stack-up and assembly complexity
    \item Defined PCB-based electrical integration strategy replacing traditional
          control cabinets; coordinated with external SMT manufacturers
    \item Implemented Siemens S7-1200 PLC for process control with custom PCB-based
          I/O and signal conditioning
    \item Led industrialisation into low-volume serial production: designed
          manufacturing routes, workstation layout, and parallel assembly workflows
    \item Built a production line from scratch; led a 6-person production team;
          directed R\&D team of 7 engineers (overall headcount $\sim$16)
    \item Defined incoming inspection criteria, module-level acceptance testing
          (assembly accuracy and leak rate), and authored complete SOP set
    \item Introduced Odoo ERP to digitise production planning and purchasing
\end{itemize}

\medskip
\noindent\textbf{Chengdu Kechuang Institute of Science and Culture}, China
\hfill 03/2019 -- 06/2019 \\
\textit{Research Engineer}

\noindent Conducted comparative research on HPHT and CVD diamond fabrication
methods and associated equipment; successfully raised 2\,M CNY in seed funding for
Sichuan 330 Semiconductor Ltd.

%====================
% Patents
%====================
\section*{Selected Patents}

\begin{itemize}[noitemsep, topsep=2pt, leftmargin=*]
    \item \textbf{High-power TE--TEM microwave mode converter} (CN212162036U, 2020) ---
          low-loss mode conversion for efficient high-power microwave transmission
    \item \textbf{Waveguide sealing window} (CN212392359U, 2021) ---
          vacuum-compatible dielectric window for high-power microwave systems
    \item \textbf{Microwave sliding short with choking structure} (CN212162044U, 2020) ---
          adjustable impedance tuning with reduced RF leakage
    \item \textbf{Large-area compact MPCVD device} (CN212560433U, 2021) ---
          modular architecture for stable large-area plasma generation
\end{itemize}

%====================
% Languages
%====================
\section*{Languages}

Chinese (native) \enspace $\cdot$ \enspace English (fluent, professional working level)

\end{document}
